% !TEX root = ../nc-quadric-surface.tex
\subsection{二型がともに階数三となる非標準表示}
\label{sec:rank-three-example}

\begin{proposition}
\label{prop:nonstandard-rank-three-example}
標数零の代数閉体上に、二つの MCM 型の行列成分がともに三次元を張る
non-standard な $\HH$ の適合表示が存在する。
この表示と定理~\ref{thm:normal-helix} のいずれの初期型を用いても、
得られる四周期表示は Van den Bergh の cubic 型にならない。
\end{proposition}

\begin{proof}
$k$ の元 $\mathrm i$ を $\mathrm i^2=-1$ となるように選び、
\[
 A=k[u,v]\otimes_k k\langle x,y\rangle/(yx+xy),\qquad
 g=\tfrac14(x^2+y^2)-uv,\qquad B=A/(g)
\]
とする。$A$ は四変数の skew polynomial ring であり、
命題~\ref{prop:commutative-twist-covers} で用いた skew Koszul 複体と同様に
型 $(1,4,6,4,1)$ である。$g$ は中央正則である。
$B$ は $k[u,v,x]$ 上で $y$ に関する Ore 拡大を取り、
モニックな関係 $y^2=4uv-x^2$ を課した環なので、
左 $k[u,v,x]$ 加群として基底 $1,y$ を持つ。
特に $u$ は正則である。$u$ を可逆にすると $v$ を消去でき、
\[
 B[u^{-1}]\simeq k[u,u^{-1}]\otimes_k k\langle x,y\rangle/(yx+xy)
\]
となるため、$B$ は domain である。

滑らかさは $B$ の Koszul 双対から判定できる。
$U,V,X,Y$ を $u,v,x,y$ の双対基底とすると、二次関係の直交補空間は
\[
\begin{gathered}
 U^2=V^2=0,\quad X^2=Y^2=w,\quad UV+VU=-4w,\quad XY=YX,\\
 UX=-XU,\quad UY=-YU,\quad VX=-XV,\quad VY=-YV
\end{gathered}
\]
を与える。ここで $w=X^2$ は中央であり、$B^!/(w)=A^!$ となる
正則元の代表である \cite[Lemma 4.11(2)]{MU2022}。$P=(U+V)/(2\mathrm i)$、$Q=(U-V)/2$ と置けば
$P^2=Q^2=w$、$PQ=-QP$ である。
$C(B)=B^![w^{-1}]_0$ は $P X^{-1},Q X^{-1},Y X^{-1}$ によって生成される。
最初の二つはそれぞれ平方が $-1$ で互いに反交換し、最後は中央で平方が $1$ である。
従って $M_2(k)\times M_2(k)$ から $C(B)$ への全射を得る。
\cite[Lemma 4.13(1)]{MU2022} により $C(B)$ は八次元なのでこの全射は同型であり、
同論文 Theorem 5.5 によって $\qgr B$ は滑らかである。

$B$ は正則二次切断なので次元三・parameter 二の AS Gorenstein 環である
\cite[Lemma 2.4]{MU2022}。canonical bimodule を与える右作用の twist を計算する。
\cite[arXiv v2, Example 5.5]{RRZ2014} により、$A$ の Nakayama 自己同型は
$u,v$ を固定し $x,y$ に符号 $-1$ を掛ける。
これは $g$ を固定するため同論文 Lemma 1.5 を適用でき、
$B$ の Nakayama 自己同型もその誘導 $\mu$ となる。
$\mu^{-1}=\mu$ なので、\cite[Example 3.6(2)]{MU2021} の
右 canonical bimodule の規約では
$S_{\qgr B}^{-1}[2]\simeq(-)_\mu(2)$ となる。

ここで graded 自己同型
\[
 \sigma(u)=-u,\qquad \sigma(v)=v,\qquad
 \sigma(x)=\mathrm i x,\qquad \sigma(y)=\mathrm i y
\]
を取る。$\sigma(g)=-g$ なので $B$ に誘導され、
その二乗は $\mu$ である。命題~\ref{prop:central-twist-obstruction} の証明中の
作用比較を用いると、Zhang 同値 $F:\qgr B\to\qgr B^\sigma$ は $\pi B$ を
$\pi B^\sigma$ に送り、新しい次数シフトを旧圏では $(-)_\sigma(1)$ として実現する。
その二乗は $(-)_\mu(2)$ なので、新しい次数シフトは反 Serre の平方根である。
$A^\sigma$ は型 $(1,4,6,4,1)$ を保ち \cite[Lemma 4.1]{MUW2026}、
$g$ はその中で左右正則正規であり、$B^\sigma=A^\sigma/(g)$ である。

$\pi B^\sigma$ に対する切断と豊富性を確認する。
正規正則切断の AS Gorenstein 性と深さ三から、
全次数での飽和性、$\operatorname{Ext}^1(\pi B^\sigma,\pi B^\sigma)=0$ が従う。
すなわち次数付き有限長加群 $L$ について $\operatorname{Ext}^j(L,B^\sigma)=0$
($j<3$) を用い、自由加群の打切りと Hom の極限を取る
\cite[Lemma 5.9(1)--(2)]{MU2021}。
$\operatorname{Ext}^2$ は Serre 双対性により負次数の切断の双対なので零である。
高次も大域次元二により零であり、$\End(\pi B^\sigma)=k$ と合わせて exceptionality を得る。
任意の対象は有限個の $\pi B^\sigma(-n)$ の直和の商であり、
その指数は打切りにより任意に大きく取れる。
全射の核 $K$ について $\operatorname{Ext}^1(\pi B^\sigma(-n),K)=0$
($n\gg0$) を得るには、$K$ の自由被覆を取って、
その核への $\operatorname{Ext}^2$ に埋め込む。
その核の自由被覆と大域次元二から、最後の群は
自由加群の像への $\operatorname{Ext}^2$ の商となり、Serre 双対性と
負次数の切断の消滅により零である。
従って $(\pi B^\sigma,(1))$ は ample であり、その切断環は $B^\sigma$ である。

Gabriel--Krull 条件は同値で不変なので $\qgr B$ で確認する。
命題~\ref{prop:sklyanin-h2} と同じく $A$ 上の有限自由分解により
Hilbert 関数は最終的に多項式となる。定数なら有限長、次数が高々一なら
$\cC_1$ に入ることは \eqref{eq:sklyanin-polynomial-filtration} と同じ議論で従う。
一方 $B/uB$ 上で $v$ は正則であり、$B/(u,v)$ は無限次元である。
従って $v$ の冪による列は $\pi(B/uB)$ の真の下降列となり、
この対象は有限長でない。さらに $u$ の冪による列の各相続商は
$\pi(B/uB)$ のシフトであるから、$\pi B\notin\cC_1$ である。
$\pi B$ の任意の非零部分対象は、非零斉次元 $f\in B$ に対する
$\pi(fB)\simeq\pi B(-\deg f)$ を含む。
$B/fB$ の Hilbert 多項式は次数が高々一なので、非零部分対象による商は
$\cC_1$ に入る。よって $\pi B$ は rank one torsion-free である。

最後に $c=(x+y)/2$、$d=(x-y)/2$ とし、$r=c,d$ に対し
\[
 M_r=\begin{pmatrix}r&u\\ -v&-r\end{pmatrix}
\]
と置く。$c^2=d^2=(x^2+y^2)/4$ により $M_r^2=gI$ であり、
これらは reduced な線形行列分解である。
成分の張る空間はそれぞれ $\langle u,v,c\rangle$ と
$\langle u,v,d\rangle$ である。これらは異なり、
行と列の定数基底変換で保たれるので、二つの余核は同型でない。
各分解では二つの連続行列が同じなので、$\Omega(1)$ は対応する余核を固定する。
この二つが全 MCM 型であることは $C(B)\simeq M_2(k)\times M_2(k)$ と
\cite[Lemma 4.13(4)--(5)]{MU2022} による。
実際、余核は二元生成で、表示行列は非零の最小関係を与えるため自由ではない。
非自由直和因子は少なくとも二元で生成されるので、各余核は直既約である。

$D=\operatorname{diag}(\mathrm i,1)$ とすると
$\sigma(M_r)=\mathrm i D M_r D^{-1}$ なので、右作用の $\sigma$ による引戻しは
両方の型を固定する。
Zhang 同値は極小自由被覆とその核を移送し、
$F(M)(1)\simeq F(M(1)_\sigma)$ であるから、
$B^\sigma$ でも $\Omega(1)$ は両方の型を固定する。
最小表示行列の成分は twist 後も同じ次元の空間を張る。
実際、$F(B(-1))\simeq B^\sigma(-1)$ を $m\mapsto\sigma^{-1}(m)$ で同定すると、
次数一成分が $a$ の左乗法は $a\star(-)$ となる。
従って $B^\sigma$ は non-standard であり、双方の成分空間が三次元である。
系~\ref{cor:normal-cubic-defect} により、どちらの初期型を選んでも
cubic 型は得られない。
\end{proof}

任意の別表示を含めた cubic 型の排除は、
\ref{thm:cubic-coverage} で証明する。

\begin{corollary}[この skew polynomial 環の中央切断の許容域]
\label{cor:skew-rank-three-central-locus}
上の $A$ の非零中央二次元 $g$ を
\[
 g=a x^2+b y^2+q(u,v)
\]
と書く。$A/(g)$ が domain かつ滑らかとなるための必要十分条件は
$ab\ne0$ かつ二次形式 $q$ が非退化となることである。
この領域では全ての中央切断およびその Zhang twist は
同じ pointed exact 同値類を与える。この類は $\HH$ に属し、
定理~\ref{thm:cubic-coverage} の cubic 表示を持たない類である。
\end{corollary}

\begin{proof}
次数二の中心は表示された五つの二次元で張られる。
これは skew polynomial の標準単項式を用い、
$x,y$ と交換する元の $x,y$ の次数がともに偶数となることから従う。
$ab\ne0$ かつ $q$ 非退化なら、$x,y$ のスカラー変更と
$u,v$ の線形基底変更で命題~\ref{prop:nonstandard-rank-three-example} の
$g$ に移る。この変更は $A$ の自己同型である。
Zhang 同値は構造対象を保存するので、この条件が十分である。

必要性を示す。$a=b=0$ なら非零の $q$ は一次元の積に分解し、
二因子は商で非零なので商は domain でない。
残りでは $x,y$ を交換して $a\ne0$ としてよい。
$B=A/(g)$ とし、$X,Y$ を $x,y$ の双対基底とする。
$B^!$ において $w=X^2$ は中央正則元の代表であり、
$XY=YX$、$Y^2=(b/a)w$ となる。
$u,v$ の双対に由来する二つの生成元は $X,Y$ と反交換し、
その平方と相互の反交換子は $w$ に $q$ の係数を掛けたものである。
これらはいずれも $B$ の二次関係空間との直交性を評価して得られる。
従って $C(B)=B^![w^{-1}]_0$ は
$Y X^{-1}$ と、$u,v$ の双対元に $X^{-1}$ を掛けた二元で生成される。

$b=0$ なら $Y X^{-1}$ は非零の中央平方零元なので $C(B)$ は
半単純でない。非零性は $w$ の正則性による。
また $q$ が退化しているなら、$u,v$ の双対側でその二次形式の根基から
非零元 $L$ を取れる。$L$ は平方が零であり、他の双対二次形式の
生成元との反交換子も零である。
よって $L X^{-1}$ は非零の平方零元で、
他の生成元と交換または反交換する。従ってそれが生成する
両側イデアルは非零で平方零となる。
いずれの場合も \cite[Theorem 5.5]{MU2022} により $\qgr B$ は滑らかでない。
\end{proof}

\begin{theorem}[四変数符号型の対角切断]
\label{thm:sign-quadric-diagonal}
$\epsilon_{ij}\in\{1,-1\}$ とし、
$A_\epsilon=k\langle x_1,\dots,x_4\rangle/
(x_jx_i-\epsilon_{ij}x_ix_j)_{i<j}$ とする。
非零中央元 $g=\sum_i a_i x_i^2$ による切断を考える。
$A_\epsilon/(g)$ が滑らかな domain となるのは、全ての $a_i$ が非零であり、
\[
 \epsilon_{ij}\epsilon_{ik}\epsilon_{jk}=-1
 \qquad(i<j<k)
\]
となる三つ組の個数が零または二の場合にちょうど限る。
前者は可換二次曲面の標識類、後者は
命題~\ref{prop:nonstandard-rank-three-example} の標識類を与える。
この二類はいずれも $\HH$ に属し、pointed exact 同値でない。
\end{theorem}

\begin{proof}
$a_j\ne0$ を一つ選ぶ。
二次双対において $w=X_j^2$ は中央正則で、
$X_i^2=(a_i/a_j)w$、$X_iX_l=-\epsilon_{il}X_lX_i$ ($i<l$) である。
$C(B)$ は $X_iX_j^{-1}$ ($i\ne j$) で生成され、
これらは相互に符号付きで交換する。
$a_i=0$ なら $X_iX_j^{-1}$ は零でない平方零元であり、
これが生成する両側イデアルも平方零である。
従って滑らかさに必要な半単純性が失われる
\cite[Theorem 5.5]{MU2022}。

以後 $a_i\ne0$ とする。対角符号自己同型 $x_i\mapsto s_i x_i$ による
Zhang twist は $\epsilon_{ij}$ を $s_i s_j\epsilon_{ij}$ に替える。
$g$ の新しい各平方に対する係数は非零のままなので、
変数のスカラー変更で正規化できる。
$g$ は自己同型で固定され、商の twist は構造対象を保つ。
また domain 性は斉次元の積に関してこの twist で不変である。

この符号変更で特定の頂点と残りの全ての間の符号を正にできる。
その後に残る三辺の符号は、その頂点を含む三つ組の積で決まる。
従って三つ組の四つの符号が各ラベルで一致するなら
二つの符号行列は twist で移り合う。
四符号の積は一であり、頂点の置換は三つ組の置換を全て実現する。
従って負の三つ組の個数は零、二、四のいずれかであり、
それぞれ全辺正、一辺だけ負、全辺負を代表として取れる。

零の場合の中央切断は滑らかな可換二次曲面である。
二の場合は可換変数二つと反交換する変数二つを持つ
前の系の非退化切断である。
四の場合には $\sum_i a_i x_i^2=(\sum_i \sqrt{a_i}x_i)^2$ となる。
一次の元は商で非零であり、平方が零なので domain ではない。
二つの残る標識類の非同値は、
可換類が cubic 表示を持ち、後者が持たないこと
（定理~\ref{thm:cubic-coverage}）から従う。
\end{proof}
